\documentclass[11pt]{article}

\usepackage[a4paper,margin=0.78in]{geometry}
\usepackage[T1]{fontenc}
\usepackage[utf8]{inputenc}
\usepackage{lmodern}
\usepackage{amsmath,amssymb,amsthm,mathtools,bm}
\usepackage{enumitem,booktabs,array,tabularx,microtype,xcolor}
\usepackage[numbers,sort&compress]{natbib}
\usepackage[colorlinks=true,linkcolor=blue!65!black,
  citecolor=blue!65!black,urlcolor=blue!70!black]{hyperref}
\usepackage[nameinlink,noabbrev]{cleveref}
\setlength{\emergencystretch}{3em}

\hypersetup{
  pdftitle={TPC-MVP4: A Source-Locked Route Decision after TPC-133--141},
  pdfauthor={Liang Wang},
  pdfsubject={Frontier totalization, arithmetic shadows, and the strict endpoint},
  pdfkeywords={fixed shift, route decision, source-locked snapshot, arithmetic frontier, endpoint}
}

\newtheorem{theorem}{Theorem}[section]
\newtheorem{proposition}[theorem]{Proposition}
\newtheorem{corollary}[theorem]{Corollary}
\theoremstyle{definition}
\newtheorem{definition}[theorem]{Definition}
\theoremstyle{remark}
\newtheorem{remark}[theorem]{Remark}

\newcommand{\GO}{\textnormal{\textsc{go}}}
\newcommand{\AINF}{\textnormal{\textsc{architecture-infeasible}}}
\newcommand{\REROUTE}{\textnormal{\textsc{reroute}}}
\newcommand{\STOPR}{\textnormal{\textsc{stop-route}}}
\newcommand{\NT}{\textnormal{\textsc{not-testable}}}
\newcommand{\AF}{\textnormal{\textsc{arithmetic-frontier}}}
\newcommand{\OPEN}{\textnormal{\textsc{open}}}
\newcommand{\INVALID}{\textnormal{\textsc{invalid-snapshot}}}
\newcommand{\PROVED}{\textnormal{\textsc{proved}}}
\newcommand{\COND}{\textnormal{\textsc{conditional}}}
\newcommand{\REFUTED}{\textnormal{\textsc{refuted}}}
\newcommand{\Lam}{\Lambda_{\mathrm{phys}}}
\newcommand{\Lzero}{\mathrm{L0}}
\newcommand{\Lone}{\mathrm{L1}}
\newcommand{\Ltwo}{\mathrm{L2}}
\newcolumntype{Y}{>{\raggedright\arraybackslash}X}
\newcolumntype{P}[1]{>{\raggedright\arraybackslash}p{#1}}

\title{\textbf{TPC-MVP4: A Source-Locked Route Decision}\\
\textbf{after TPC-133--141: Frontier Totalization,}\\
\textbf{Arithmetic Shadows, and the Strict Endpoint}}
\author{Liang Wang\\
School of Artificial Intelligence and Automation,\\
Huazhong University of Science and Technology, Wuhan 430074, P.R. China\\
\texttt{wangliang.f@gmail.com}}
\date{July 2026}

\begin{document}
\maketitle

\begin{abstract}
TPC-MVP3 reported the requirement for a complete native archive as
the first missing object in a proposed fixed-shift prime-pair route.
TPC-133--136 now make the native entrance executable and follow it
through an exact declared dyadic cut.  They reconnect the
eligible-prefix, eligible-tail and frontier terminal classes to the
original packet, prove the eligible prefix soft, and give a formal
support-envelope path-matrix witness that the eligible-only TPC-17/18
compiler is coefficientwise incomplete.  The witness coefficient
itself may vanish.
This cut completeness is not full-carrier totalization.
TPC-137--140 separately prove a fixed-data logarithmic arithmetic
shadow, a small-polylog affine power-of-log estimate outside a small
logarithmic-density exceptional set, and two scoped non-transfer
firewalls, then formulate the remaining actual-family and selector
estimates needed by the physical determinant-two packet.

We read these results only through the source-locked TPC-141
manifest.  Hash equality is only a drift/integrity check; it does not
prove a source theorem or generator.  Invalid hashes, scopes,
carriers, normalizations, terminal covers, occurrence registries or
endpoint records return no mathematical verdict.  A valid snapshot is
classified by seven
ordered outcomes:
\[
\begin{gathered}
\textsc{go},\quad \textsc{architecture-infeasible},\quad
\textsc{reroute},\\
\textsc{stop-route},\quad \textsc{not-testable},\quad
\textsc{arithmetic-frontier},\quad \textsc{open}.
\end{gathered}
\]
The arithmetic-frontier label requires all selected-route structural
nodes and the strict physical \(1/400\) endpoint to be proved
independently of the unresolved arithmetic targets.  Fixed,
logarithmic, almost-scale, conditional and scope-mismatched results
cannot trigger it.

The small-polylog result is a restricted external positive arithmetic
shadow, not positive \(\mathrm{L2}\): actual CRT-family containment,
squarefree/periodic reassembly, local exceptional-set control on every
terminal window and deterministic prefixes are not proved.  The
frozen snapshot is valid, but its first missing node is now
\(\mathsf{H1.frontier\_totalization}\).  The frontier has neither
complete downstream physical maps nor a complete original-scale
\(o(X)\) theorem; H6--H9 remain untestable; and no positive
\(\mathrm{L2}\) estimate occurs.  Hence
\[
 \boxed{\operatorname{Verdict}_{133:141}
 =\textsc{not-testable}.}
\]
This is genuine movement of the first-failure boundary, not
\textsc{go}, not the arithmetic frontier, and not a negative theorem
about the twin-prime conjecture.
\end{abstract}

\noindent\textbf{Keywords:} fixed shift; source-locked snapshot;
route decision; unsupported carrier; arithmetic frontier; parity
barrier.

\medskip
\noindent\fcolorbox{red!55!black}{red!3}{%
\begin{minipage}{0.96\linewidth}
\textbf{Snapshot boundary.}
This decision is dated 27 July 2026 and applies only to the declared
TPC-MVP1 architecture, the selected frontier-augmented signed route,
and the exact TPC-141 source manifest.  A changed source hash, route,
scope, carrier, normalization or occurrence registry creates a new
snapshot; an unchanged hash supplies integrity, not mathematical
validity.  No conclusion here asserts a prime-pair theorem or an
impossibility result for twin primes.
\end{minipage}}

\section{What changed after MVP3}

MVP1 resolved a conditional proof program into H1--H9
\citep{WangMVP1}.  MVP2 separated proved interfaces from missing
growing inputs \citep{WangTPC120}.  MVP3 added a strict
arithmetic-frontier label and reported
\(\mathsf{H1.archive}\) as the first missing record
\citep{WangTPC132}.  TPC-131 supplied the general composition formula
\[
 |B_{h_0,\delta}(X)|
 \le o(X)+X^{1-\sigma_{\rm raw}+\Lambda_{\rm phys}+o(1)}
\label{eq:synthesis}
\]
under a complete proof-carrying archive and occurrence registry
\citep{WangTPC131}.

TPC-133--136 provide more than another abstract schema.  They
enumerate the native entrance in an effective support envelope,
create a coefficientwise conservative dyadic path archive, partition
the TPC-17/18 eligible domain from its frontier, and reconnect the
three declared cut-terminal classes exactly
\citep{WangTPC133,WangTPC134,WangTPC135,WangTPC136}.  Thus the old
generic archive absence has been replaced by a precise downstream
totality defect.

The arithmetic branch also becomes sharper but not stronger in the
required sense.  TPC-137 proves a frozen fixed-affine logarithmic
closure based on the formal two-point theorem of Tao
\citep{Tao2016,WangTPC137}.  TPC-138 records the quantitative
shift-one scope of Helfgott--Radziwi{\l}{\l} and proves that the
direct reparameterization is blocked on the active odd
determinant-two carrier \citep{HelfgottRadziwill2021,WangTPC138}.
TPC-138/139 also record the Tao--Ter\"av\"ainen small-polylog affine
almost-scale theorem as a restricted positive shadow.  TPC-139 and
TPC-140 then expose, rather than assume, the remaining actual-family
containment, squarefree/periodic reassembly, local exceptional-window
and actual-prefix selector gates \citep{WangTPC139,WangTPC140}.  In
particular, a cumulative logarithmic-density estimate is not promoted
to the same normalized estimate on every terminal window.

\section{A frozen source-locked snapshot}

\begin{definition}[Snapshot validity]
\label{def:valid}
An MVP4 snapshot is valid only if:
\begin{enumerate}[label=(\roman*),leftmargin=2.2em]
\item the TPC-141 manifest hash and every TPC-133--140 source bundle
  hash agree with the frozen files;
\item node, export, import, scope, carrier and normalization tables
  are internally consistent;
\item the proof DAG is acyclic and its declared topological order
  exactly covers the nodes;
\item the three cut terminal roles and four residual namespaces are
  exact covers;
\item the occurrence registry has unique identifiers, acyclic
  nonoverlapping joint replacements and no missing primitive leaf;
\item the selected route and every stop record belong to the declared
  route universe; and
\item the endpoint uses exact rational direction: upper certificates
  pass, lower certificates stop, and unknown values remain unknown.
\end{enumerate}
An invalid snapshot returns \(\INVALID\), which is a data-validation
outcome rather than a route verdict.
\end{definition}

\begin{remark}[Hash scope]
A content hash detects source drift.  It does not prove a theorem,
the completeness of a generator, or the uniqueness of a tuple key.
Those claims require the source proofs and raw tuple validation.
\end{remark}

\section{Statuses, evidence and active requirements}

Node status and program evidence remain orthogonal:
\[
\PROVED,\quad \COND,\quad \OPEN,\quad \NT,\quad \REFUTED.
\]
An \(\Lzero\) record is a finite or abstract identity.  An
\(\Lone\) record is attached to the literal fixed-\(h_0\) carrier and
normalization.  A positive \(\Ltwo\) record must be a uniform growing
power estimate on the complete actual coefficient family with all
masks, weights, phases, origins, prefixes and residual sectors.

Accordingly:
\begin{itemize}[leftmargin=2em]
\item fixed-data or fixed-periodic logarithmic cancellation is an
  arithmetic shadow, not positive \(\Ltwo\);
\item a small-polylog affine power-of-log estimate outside an
  exceptional scale set is a restricted positive arithmetic shadow,
  not actual-family containment, a local terminal-window estimate, or
  deterministic-prefix control;
\item a shift-one estimate is not an estimate on a scope-mismatched
  determinant-two progression;
\item a non-effective slow diagonal is not the prescribed physical
  schedule;
\item almost-all-scale cancellation is not all-prefix control; and
\item a conditional power-ledger formula is not a certified positive
  exponent.
\end{itemize}

Route requirements are an explicit selected AND/OR graph.  A stopped
positive or eligible-only sufficient subroute need not be proved on a
selected signed alternative.  Conversely, one arithmetic theorem
cannot cover two packet classes unless a joint exact-cover theorem
says so.

\section{Eight total audit outcomes}

The outer audit has one validation outcome and seven mathematical
route verdicts.

\begin{definition}[Ordered valid-snapshot verdicts]
\label{def:verdicts}
After \cref{def:valid} passes, test:
\begin{enumerate}[label=(\roman*),leftmargin=2.2em]
\item \(\GO\): every active required gate is proved in one compatible
  scope and the actual endpoint has strict positive slack.
\item \(\AINF\): a theorem-backed exact cover of the declared route
  universe stops every route.
\item \(\REROUTE\): the selected route is stopped and a distinct
  typed alternative with a fresh registry has been selected.
\item \(\STOPR\): the selected route has a complete actual-carrier
  stop certificate and no typed alternative is selected.
\item \(\NT\): an active required artifact or type is absent.
\item \(\AF\): every active structural node and the strict physical
  endpoint are proved, while every unresolved active node is
  status-\(\OPEN\), scope-compatible and explicitly a positive
  \(\Ltwo\) target.
\item \(\OPEN\): every remaining valid case.
\end{enumerate}
\end{definition}

\begin{theorem}[Total deterministic classifier]
\label{thm:classifier}
Manifest validation followed by the displayed precedence returns
exactly one of
\[
\begin{gathered}
\INVALID,\quad \GO,\quad \AINF,\quad \REROUTE,\\
\STOPR,\quad \NT,\quad \AF,\quad \OPEN.
\end{gathered}
\]
\end{theorem}

\begin{proof}
Validation either fails, giving the first outcome, or supplies a
well-typed route universe, selected route, gate table and endpoint.
On a valid snapshot, \cref{def:verdicts} is an ordered list ending in
a catch-all.  The first applicable predicate is unique by precedence.
The order deliberately separates global stop cover, selected-route
stop, missing data and the arithmetic frontier.
\end{proof}

\begin{remark}[Conditional synthesis is not a verdict]
The existence of \eqref{eq:synthesis} is metadata: it tells us what
would follow from a complete certificate.  It is neither a weakened
\(\GO\) nor evidence that the remaining inputs are technical.
\end{remark}

\section{The strict endpoint and stop directions}

For the MVP1 contract the physical target is
\[
 \Lambda_{\rm phys}<\frac1{400}
\]
on the amplitude scale.  The TPC-141 occurrence registry keeps
energy-to-amplitude conversion, joint replacement and packet census
inside this ledger.

\begin{proposition}[Upper failure is not a lower obstruction]
\label{prop:directions}
For one fully specified route cell:
\begin{enumerate}[label=(\alph*),leftmargin=2em]
\item a complete theorem-backed upper certificate
  \(U<1/400\) passes the physical endpoint;
\item a complete theorem-backed lower or exact dual certificate
  \(L\ge1/400\) stops the route cell; and
\item a conservative upper certificate \(U\ge1/400\), without such a
  lower certificate, merely fails to prove the endpoint and does not
  establish \(\STOPR\).
\end{enumerate}
Unknown tokens give \(\NT\), not zero.  H5 determinant reserve cannot
be subtracted from \(U\).
\end{proposition}

This distinction is necessary because a poor upper estimate may be
improvable, while a lower certificate excludes every allowed
realization of the selected cell.

\section{H9 must be arithmetically independent}

The physical occurrence registry is a structural certificate.  In
the MVP4 proof DAG,
\[
\mathsf{H9.physical\_registry}
\longleftarrow
\{\mathsf{H1.actual\_carrier},\mathsf{H6.cover},
\mathsf{H7.fixed\mbox{-}h_0},\mathsf{H8.reconnection}\},
\]
with no direct or indirect dependency path from H2, H3, H4 or H5.

\begin{proposition}[Noncircular arithmetic-frontier test]
\label{prop:h9}
If H9 is proved using a target H2 or H3 saving, a complete H4
arithmetic tail estimate, or H5 determinant reserve as a physical
token, the snapshot cannot be labeled \(\AF\).
\end{proposition}

\begin{proof}
The arithmetic-frontier predicate is intended to assert that all
nonarithmetic interfaces are already closed before the remaining
arithmetic targets are attacked.  Using one of those targets to prove
the physical registry assumes a downstream premise and makes that
predicate circular.  TPC-141 therefore rejects such an H9 parent
edge.
\end{proof}

\section{Projection of TPC-133--141}

\begin{table}[htbp]
\centering
\small
\caption{MVP4 gate projection in the frozen snapshot.}
\label{tab:gates}
\begin{tabularx}{\textwidth}{@{}P{0.07\textwidth}P{0.18\textwidth}
P{0.18\textwidth}Y@{}}
\toprule
Gate & Status & Evidence target & First unresolved content\\
\midrule
H1 & \(\NT\) & \(\Lone\) &
Frontier totalization after the proved cut archive.\\
H2 & \(\OPEN\) & positive \(\Ltwo\) &
Signed replacement for the unresolved resonance sector.\\
H3 & \(\OPEN\) & positive \(\Ltwo\) &
Actual CRT containment, squarefree/periodic reassembly,
local exceptional-window control, selector/all-prefix and positive
fixed-\(X\)-power raw-saving theorem.\\
H4 & \(\OPEN\) & positive \(\Ltwo\) &
Complete original-scale physical high/ultra/boundary return.\\
H5 & \(\NT\) & positive \(\Ltwo\) &
Actual determinant/zero pair on totalized fibers.\\
H6 & \(\NT\) & \(\Lone\) &
Physical cover including every frontier leaf.\\
H7 & \(\NT\) & \(\Lone\) &
Total fixed-\(h_0\) map on every terminal leaf.\\
H8 & \(\NT\) & \(\Lone\) &
Final hard-packet intertwining beyond the cut.\\
H9 & \(\NT\) & \(\Lone\) &
Complete occurrence registry and strict physical primal.\\
\bottomrule
\end{tabularx}
\end{table}

The selected route is
\nolinkurl{signed_native_frontier_augmented}.  TPC-135 stops the
\nolinkurl{eligible_only_tpc17_18_compiler}, and TPC-138 stops
the direct \nolinkurl{quantitative_shift1_reparameterization}.  Neither
witness stops the selected frontier-augmented route, so the declared
route universe has no complete stop cover.

\section{First missing and current verdict}

The TPC-141 topological scan returns
\[
\boxed{
\begin{aligned}
\mathrm{node}&=\mathsf{H1.frontier\_totalization},\\
\mathrm{status}&=\NT,\\
\mathrm{artifact}&=
\text{total downstream frontier maps or a complete frontier theorem}.
\end{aligned}}
\]

\begin{theorem}[MVP4 route decision]
\label{thm:decision}
For the source-locked TPC-133--141 snapshot dated 27 July 2026,
\[
 \boxed{\operatorname{Verdict}_{133:141}=\NT.}
\]
The result is neither \(\GO\), nor \(\AF\), nor \(\AINF\).
\end{theorem}

\begin{proof}
The source manifest passes validation, so \(\INVALID\) does not
apply.  The selected route is not stopped and the route universe is
not completely covered by stop witnesses, so the three stop outcomes
do not apply.  The required frontier-totalization node is absent;
hence \(\NT\) occurs before \(\AF\) and \(\OPEN\).  Independently,
H6--H9 and the strict endpoint are incomplete, and no positive
\(\Ltwo\) theorem is present, ruling out both \(\GO\) and \(\AF\).
\end{proof}

\begin{corollary}[Precise progress statement]
The first-failure pointer has moved from a generic missing native
archive to the first unsupported carrier after an exact cut.  This is
structural \(\Lone\) progress.  It is not a positive \(\Ltwo\)
arithmetic event.
\end{corollary}

\section{Publishable outcomes and next branch}

The batch has publishable content even though the route verdict is
\(\NT\):
\begin{itemize}[leftmargin=2em]
\item executable native enumeration and exact cut conservation;
\item a scoped formal coefficientwise obstruction to the eligible-only
  compiler on the support-envelope archive;
\item a fixed-data logarithmic arithmetic closure;
\item a restricted small-polylog affine almost-scale power-of-log
  arithmetic shadow;
\item a scoped non-transfer from quantitative shift one;
\item exceptional-scale and selector firewalls; and
\item a source-locked integration theorem with a directional
  endpoint calculus.
\end{itemize}

A bare missing file is not a negative theorem.  The publishable
content is the exact theorem that reaches and classifies the missing
frontier, together with the two scoped route obstructions.

The next minimal branch is now forced:
\begin{enumerate}[leftmargin=2em]
\item enumerate every frontier terminal from TPC-136;
\item attach \(Q_D,Q_Z,G\) and fixed-\(h_0\) maps, or prove the
  complete frontier sum \(o(X)\);
\item rebuild H6--H9 and the occurrence registry without inheriting
  old slack;
\item certify \(\Lambda_{\rm phys}<1/400\) independently of H3; and
\item only then test whether all remaining unresolved nodes are
  actual positive-\(\Ltwo\) targets.
\end{enumerate}

If the frontier totalizes and the physical endpoint passes, but H2--H5
remain the only open compatible arithmetic targets, a new frozen
snapshot may legitimately return \(\AF\).  If an actual kernel or
endpoint dual stops the frontier-augmented route, the correct output
is \(\STOPR\) or \(\REROUTE\), not a claim about twin primes.

\section{Deterministic regression and claim boundary}

The standard-library script
\texttt{experiments/tpc142\_mvp4\_route\_audit.py} reads the frozen
TPC-141 manifest and checks:
\begin{enumerate}[leftmargin=2em]
\item manifest hash, source-lock summary, DAG and first-missing
  consistency;
\item direct and indirect independence of H9 from H2--H5;
\item reachability of all seven valid-snapshot verdicts;
\item rejection of a malformed snapshot before verdict selection;
\item rejection of conditional, frozen and refuted pseudo-frontiers;
\item refusal to infer architecture infeasibility from an incomplete
  route cover;
\item refusal to turn an upper-bound failure into a stop; and
\item absence of a proved positive-\(\Ltwo\) record.
\end{enumerate}
Default mode writes a frozen snapshot and audit JSON.
\texttt{--check} performs no writes and compares both artifacts byte
for byte.

\paragraph{Established.}
Snapshot validation, ordered verdict selection, the first-missing
projection, H9 independence and endpoint direction rules are
\(\Lzero/\Lone\) audit results.

\paragraph{Not established.}
No complete frontier theorem, H1--H9 conjunction, strict physical
endpoint, positive fixed-\(h_0\) \(\Ltwo\) saving,
\(B_{h_0,\delta}=o(X)\), Hardy--Littlewood asymptotic, prime-pair
lower bound, or twin-prime theorem is proved.

\section{Conclusion}

MVP4 records real progress without changing the arithmetic claim
boundary.  The native archive has become executable through a
complete declared cut, not a totalized full carrier, and a false
eligible-only compiler has been stopped.
Known Chowla information has been placed in its exact frozen,
shift-one, or restricted small-polylog almost-scale scope.  The
remaining structural failure is now one named frontier-totalization
node, not a vague missing dictionary.

Because that node and the physical registry are incomplete, the
route is not yet at the arithmetic frontier.  The correct verdict is
\(\NT\).  This verdict is useful precisely because its source,
scope, first missing artifact, scoped negative results and permitted
next route are all frozen and machine-auditable.

\begingroup
\small
\setlength{\bibsep}{2pt plus 0.3ex}
\sloppy
\bibliographystyle{plainnat}
\bibliography{references}
\endgroup

\end{document}
